\documentclass[12pt,a4paper]{article}
\usepackage{iftex}
\ifPDFTeX
  \usepackage[T2A]{fontenc}
  \usepackage[utf8]{inputenc}
  \usepackage[english,russian]{babel}
\else
  \usepackage{fontspec}
  \setmainfont{DejaVu Serif}
  \setmonofont{DejaVu Sans Mono}
  \renewcommand{\contentsname}{Содержание}
  \renewcommand{\figurename}{Рис.}
  \renewcommand{\tablename}{Таблица}
\fi

\usepackage[a4paper,margin=25mm]{geometry}
\usepackage{amsmath,amssymb,amsthm}
\usepackage{graphicx}
\usepackage{booktabs}
\usepackage{enumitem}
\usepackage[hidelinks]{hyperref}

\newcommand{\Exp}{\mathsf{M}}
\newcommand{\Var}{\mathsf{D}}
\newcommand{\PP}{\mathsf{P}}
\newcommand{\dens}{\mathrm{p}}
\newcommand{\Normal}{\mathcal{N}}
\newcommand{\R}{\mathbb{R}}
\newcommand{\xs}[1]{x^{(#1)}}
\newcommand{\ys}[1]{y^{(#1)}}
\newcommand{\norm}[1]{\left\lVert #1 \right\rVert}
\newcommand{\abs}[1]{\left\lvert #1 \right\rvert}
\newcommand{\dd}{\,\mathrm{d}}
\newcommand{\est}[1]{\tilde{#1}}
\newcommand{\zad}[1]{\medskip\noindent\textbf{Задание #1.}\quad}
\newcommand{\prov}{\smallskip\noindent\textit{Проверка:}\ }

\theoremstyle{definition}
\newtheorem{definition}{Определение}
\newtheorem{statement}{Утверждение}
\theoremstyle{plain}
\newtheorem{lemma}{Лемма}

% Список с русскими буквенными метками
\newenvironment{ruslist}
  {\begin{itemize}[label={},leftmargin=2.2em,itemsep=2pt,topsep=3pt]}
  {\end{itemize}}


\begin{document}

\begin{center}
  {\large\bfseries Задача 5. Деревья решений}\\[1mm]
  Нейронные сети и машинное обучение\\[1mm]
  Задания 1--9
\end{center}

\section{Что нужно знать}

\subsection{Обозначения}

\begin{center}
\begin{tabular}{ll}
\toprule
$S$, $S_i$ & выборка и подвыборка $i$-го узла \\
$N$, $N_i$ & их размеры \\
$c$ (или $K$) & число классов \\
$d$ & число атрибутов (размерность вектора признаков) \\
$\PP_k = N_i^{(k)}/N_i$ & доля элементов $k$-го класса в подвыборке \\
$\mathcal{I}(S)$ & мера неоднородности выборки \\
$H(S)$ & энтропия \\
$\Delta\mathcal{I}$ & уменьшение неоднородности при разбиении \\
$\log$ & логарифм по основанию 2, единица измерения -- бит \\
\bottomrule
\end{tabular}
\end{center}

\subsection{Дерево решений}

Дерево решений классифицирует объект последовательностью вопросов. Каждому
внутреннему узлу сопоставлен классификатор (вопрос об атрибутах объекта), рёбрам
-- возможные ответы, листьям -- метки классов. Объект спускается от корня к
листу, и метка листа объявляется ответом.

Узлу $i$ отвечает подвыборка $S_i$ -- множество элементов обучающей выборки,
попадающих в этот узел. Разбиение узла делит $S_i$ на непересекающиеся
подвыборки $S_{ij}$ по числу ответов на вопрос.

Классификаторы в узлах бывают двух типов:
\begin{ruslist}
  \item[--] \emph{покоординатные}, проверяющие один атрибут (например
    $x_i \le \theta$). Дерево из таких узлов называется \emph{монотетичным};
  \item[--] \emph{не покоординатные}, использующие сразу несколько атрибутов
    (например $\sigma(w^{\top}x) > 1/2$). Дерево называется
    \emph{политетичным}.
\end{ruslist}

\subsection{Меры неоднородности}

Разбиение выбирают так, чтобы подвыборки становились однороднее, то есть
содержали преимущественно элементы одного класса. Однородность измеряют одной из
функций неоднородности:
\[
  \text{энтропия}\quad H(S_i) = -\sum_{k=1}^{c}\PP_k\log \PP_k,
  \qquad
  \text{Джини}\quad \mathcal{I}_2(S_i) = 1 - \sum_{k=1}^{c}\PP_k^2,
\]
\[
  \text{ошибка классификации}\quad \mathcal{I}_3(S_i) = 1 - \max_k \PP_k .
\]
Здесь $0\log 0 = 0$. Все три обращаются в нуль на однородной выборке и
максимальны при равных долях классов.

Качество разбиения оценивают уменьшением неоднородности:
\begin{equation}\label{eq:gain}
  \Delta\mathcal{I} = \mathcal{I}(S_i) - \sum_j \frac{N_{ij}}{N_i}\,\mathcal{I}(S_{ij})
  \;\longrightarrow\; \max .
\end{equation}
Вычитается средневзвешенная неоднородность потомков, веса -- доли попавших в них
элементов.

\subsection{Энтропия как мера информации}

Для задания 7 понадобится вероятностная запись. Пусть $Y$ -- случайная метка
класса элемента, выбранного из $S_i$ равномерно, а $Z$ -- ответ проверяемого в
узле вопроса для этого элемента. Тогда $H(S_i) = H(Y)$, а средневзвешенная
энтропия потомков есть в точности условная энтропия:
\[
  \sum_j \frac{N_{ij}}{N_i}H(S_{ij}) = \sum_j \PP(Z = j)\,H(Y \mid Z = j) = H(Y \mid Z).
\]
Поэтому \eqref{eq:gain} для энтропии -- это взаимная информация
\begin{equation}\label{eq:mi}
  \Delta H = H(Y) - H(Y\mid Z) = I(Y; Z),
\end{equation}
для которой известны свойства $I(Y;Z) = I(Z;Y) \le H(Z)$ и
$H(Z) \le \log(\text{число значений } Z)$.

\subsection{Суррогатные расщепления}

В пособии раздел про неопределённые атрибуты оставлен пустым, поэтому опишу
метод, о котором идёт речь в задании 8.

Пусть в узле выбрано основное разбиение по атрибуту $x_s$ (то, что даёт
максимум \eqref{eq:gain}), но у некоторого объекта значение $x_s$ не задано. Тогда
поступают так: среди остальных атрибутов ищут тот, разбиение по которому
максимально совпадает с основным на элементах обучающей выборки, у которых
задано и то и другое. Такое разбиение называют \emph{суррогатным} и применяют
к объекту вместо основного. Строится обычно несколько суррогатов, упорядоченных
по степени совпадения с основным.

\section{Задания}

\zad{1}
\dano Произвольное дерево решений.

\treb Доказать, что существует функционально эквивалентное ему дерево решений, в
котором последовательность классификаторов на пути от корня к листу содержит
только неповторяющиеся классификаторы.

\resh Функциональная эквивалентность означает, что деревья дают одинаковые
ответы на всех возможных объектах.

Пусть на некотором пути от корня к листу один и тот же классификатор $h$
встречается дважды: в узле $u$ и в лежащем ниже узле $v$. Все объекты,
доходящие до $v$, прошли через $u$ и свернули в конкретного потомка, то есть все
они дали на вопрос $h$ один и тот же ответ $a$. Классификатор $h$ детерминирован
и зависит только от объекта, поэтому в узле $v$ те же объекты дадут тот же ответ
$a$. Значит из $v$ достижим ровно один потомок -- тот, что отвечает ответу $a$; в
остальные не попадает ни один объект.

Заменим узел $v$ на этого единственного достижимого потомка вместе с его
поддеревом, а остальные поддеревья удалим. Ни один объект при этом не изменит
своего маршрута, значит функция, вычисляемая деревом, не изменилась. При этом
число узлов строго уменьшилось.

Повторяем операцию, пока на путях остаются повторы. Так как число узлов конечно
и на каждом шаге строго убывает, процесс закончится. Полученное дерево
эквивалентно исходному, и ни на одном пути от корня к листу классификатор не
повторяется. $\blacksquare$

\zad{2}
\dano Произвольное дерево решений, в узлах которого возможны классификаторы с
любым конечным числом ответов.

\treb Доказать, что существует эквивалентное ему двоичное дерево решений.

\resh Достаточно показать, как заменить один узел с $B > 2$ потомками, а затем
применить замену ко всем таким узлам.

Пусть узел $u$ проверяет классификатор $h$ с возможными ответами
$a_1, \dots, a_B$, и потомок $T_j$ отвечает ответу $a_j$. Заменим $u$ цепочкой из
$B-1$ двоичных узлов $u_1, \dots, u_{B-1}$, где узел $u_j$ проверяет вопрос
\[
  h_j(x) = \bigl[\,h(x) = a_j\,\bigr],
\]
то есть «верно ли, что ответ равен $a_j$». Это допустимый классификатор: он
вычислим по объекту и имеет два ответа. При ответе «да» переходим в поддерево
$T_j$, при ответе «нет» -- в узел $u_{j+1}$. Последний узел $u_{B-1}$ при ответе
«нет» ведёт в $T_B$, так как остальные варианты уже исключены.

Объект, дававший в исходном дереве ответ $a_j$, пройдёт по цепочке $j-1$ отказов
и на $j$-м вопросе свернёт в $T_j$ -- то же самое поддерево, что и раньше. Значит
функция не изменилась.

Применяя замену ко всем узлам с более чем двумя потомками (их конечное число),
получаем двоичное дерево, эквивалентное исходному. $\blacksquare$

\zad{3}
\dano Дерево с двумя уровнями: корневая вершина смежна с $B$ листьями.

\treb Найти нижнюю и верхнюю границы количества уровней в функционально
эквивалентном двоичном дереве как функции от $B$, а также количество узлов в
таком дереве.

\resh Эквивалентное двоичное дерево должно различать те же $B$ исходов, то есть
иметь $B$ листьев.

\emph{Нижняя граница.} Двоичное дерево глубины $D$ (то есть с $D$ уровнями
внутренних узлов) имеет не более $2^{D}$ листьев. Из $2^{D} \ge B$ получаем
$D \ge \lceil \log_2 B\rceil$. Всего уровней, считая уровень листьев,
\[
  L_{\min} = \lceil \log_2 B\rceil + 1 .
\]
Эта граница достигается на сбалансированном дереве, в котором вопросы делят
множество оставшихся вариантов пополам.

\emph{Верхняя граница.} Худший случай -- цепочка из построения задания 2: $B-1$
последовательных вопросов вида «ответ равен $a_j$?». Глубина такой цепочки равна
$B-1$, всего уровней
\[
  L_{\max} = B .
\]
Дерево длиннее строить бессмысленно: по заданию 1 повторяющиеся классификаторы
удаляются, а различных полезных вопросов в цепочке не больше $B-1$.

\emph{Количество узлов.} В любом двоичном дереве, где каждый внутренний узел
имеет ровно двух потомков, число внутренних узлов на единицу меньше числа
листьев. Это видно из подсчёта рёбер: каждый узел кроме корня входит ровно в одно
ребро, а каждый внутренний узел порождает два ребра, откуда
$2 \cdot \#\text{внутр} = \#\text{всего} - 1$ и $\#\text{всего} = 2B - 1$. Итак
\[
  \#\text{листьев} = B, \qquad
  \#\text{внутренних} = B - 1, \qquad
  \#\text{всего} = 2B - 1,
\]
причём это количество не зависит от формы дерева: сбалансированное и вырожденное
деревья различаются глубиной, но не числом узлов. $\blacksquare$

\zad{4}
\dano Выборка из $N$ элементов, по $N/2$ каждого из двух классов. Элемент имеет
$d$ атрибутов, каждый принимает одно из $k$ значений. Наилучшее разбиение делит
выборку на две равные части. Сравниваются монотетичное дерево (вопрос по одному
атрибуту) и политетичное (вопрос по совокупности атрибутов).

\treb (а) Сколько уровней будет в каждом дереве? (б) Какова сложность нахождения
оптимального разбиения в корне? (в) Сравнить общую вычислительную сложность
обучения.

\resh

\emph{(а) Уровни.} По условию каждое разбиение делит подвыборку пополам, причём
одинаково для обоих типов деревьев. Размер подвыборки на уровне $l$ равен
$N/2^{l-1}$, ветвление прекращается при размере 1, откуда
\[
  \frac{N}{2^{L-1}} = 1 \;\Longrightarrow\; L = \log_2 N + 1 .
\]
Число уровней одинаково у обоих деревьев: предположение о делении пополам
уравнивает их по глубине. Вся разница окажется в стоимости выбора разбиения.

\emph{(б) Сложность выбора разбиения в корне.} Оценим число кандидатов и цену
проверки одного кандидата. Проверка любого разбиения требует один раз пройти по
выборке и посчитать классовые частоты в частях, то есть $O(N)$ операций.

Монотетичное дерево. Вопрос задаётся выбором атрибута ($d$ вариантов) и
разбиением множества его $k$ значений на две группы ($2^{k-1} - 1$ вариантов).
Кандидатов
\[
  d\,(2^{k-1} - 1),
\]
и полная стоимость перебора в корне $O\bigl(d\,2^{k}N\bigr)$. Если
ограничиться вопросами вида «атрибут равен $v$», кандидатов всего $dk$, а
стоимость $O(dkN)$.

Политетичное дерево. Вопрос может опираться на все атрибуты сразу, то есть
разбивает множество всех возможных комбинаций значений, которых $k^d$. Число
разбиений этого множества на две группы равно
\[
  2^{\,k^d - 1} - 1,
\]
и полная стоимость перебора $O\bigl(2^{\,k^d}N\bigr)$.

\emph{(в) Общая сложность обучения.} На уровне $l$ находится $2^{l-1}$ узлов, в
каждом по $N/2^{l-1}$ элементов, суммарно на уровень приходится $N$ элементов.
Значит стоимость одного уровня равна стоимости обработки $N$ элементов при том же
числе кандидатов, а всего уровней $\log_2 N + 1$:
\[
  T_{\text{мон}} = O\bigl(d\,2^{k}\,N\log N\bigr),
  \qquad
  T_{\text{пол}} = O\bigl(2^{\,k^d}\,N\log N\bigr) .
\]
Отношение
\[
  \frac{T_{\text{пол}}}{T_{\text{мон}}} = O\!\left(\frac{2^{\,k^d}}{d\,2^{k}}\right)
\]
растёт дважды экспоненциально по $d$. Уже при $k = 2$ и $d = 6$ показатель равен
$2^{64}$ против $d\cdot 2^2 = 24$. Именно поэтому на практике полный перебор
политетичных разбиений не используется: вместо него в узле обучают
параметрический классификатор (логистическую регрессию), сводя перебор к
оптимизации. $\blacksquare$

\zad{5}
\dano Монотетичное дерево решений, обучаемое с использованием энтропии.
$c$ классов, сбалансированная выборка из $n$ элементов ($n/c$ на класс),
элементы $d$-мерные.

\treb Оценить вычислительную сложность обучения по пунктам: (а) число разбиений
в корне; (б) сложность нахождения разбиения, минимизирующего энтропию;
(в) сложность этого разбиения в корне; (г) число уровней при делении подвыборок
пополам до листьев из одного элемента; (д) число классов на каждом уровне;
(е) сложность построения дерева до $L$-го уровня; (ж) время обучения и
прогнозирования при $n = 2^{10}$, $d = 6$, $c = 16$ и $10^{-10}$ с на операцию.

\resh

\emph{(а) Число разбиений в корне.} Покоординатный вопрос имеет вид
$x_i \le \theta$. Различные $\theta$ дают различные разбиения только когда порог
проходит между соседними значениями $i$-й координаты, а таких промежутков не
больше $n-1$. Умножая на число атрибутов,
\[
  \#\text{разбиений} \le d\,(n-1) = O(dn).
\]

\emph{(б) Сложность выбора наилучшего разбиения.} Разберём один атрибут.
Отсортируем по нему элементы: $O(n\log n)$. Далее ведём порог слева направо. При
переносе одного элемента через порог меняются два классовых счётчика, а
пересчёт энтропий обеих частей стоит $O(c)$. Всего $n-1$ положений порога,
поэтому проход стоит $O(nc)$. Для одного атрибута $O(n\log n + nc)$, для всех
\[
  O\bigl(dn(\log n + c)\bigr).
\]
Наивная реализация, пересчитывающая частоты заново для каждого порога, стоила бы
$O(dn^2c)$; выигрыш даёт именно инкрементальный пересчёт.

\emph{(в) В корне.} Подставляем размер выборки $n$: сложность равна
$O\bigl(dn(\log n + c)\bigr)$.

\emph{(г) Число уровней.} Размер подвыборки на уровне $l$ равен $n/2^{l-1}$,
лист достигается при размере 1:
\[
  \frac{n}{2^{L-1}} = 1 \;\Longrightarrow\; L = \log_2 n + 1 .
\]

\emph{(д) Число классов на уровне.} В узле уровня $l$ содержится $n_l = n/2^{l-1}$
элементов, значит различных классов там не больше $n_l$; с другой стороны их не
больше $c$. Отсюда
\[
  c_l = \min\Bigl(c,\ \frac{n}{2^{l-1}}\Bigr).
\]
Проверим по условию: на первом уровне $c_1 = \min(c, n) = c$ -- все классы;
в узлах, смежных листьям, подвыборка состоит из двух элементов, и
$c_l = \min(c, 2) = 2$ -- два класса. Пока $n_l > c$, число классов остаётся
равным $c$, а затем убывает вдвое на каждом уровне.

\emph{(е) Сложность до $L$-го уровня.} На уровне $l$ находится $2^{l-1}$ узлов, в
каждом $n_l = n/2^{l-1}$ элементов. Стоимость уровня
\[
  2^{l-1}\cdot O\bigl(d\,n_l(\log n_l + c_l)\bigr)
  = O\bigl(d\,n\,(\log n_l + c_l)\bigr),
\]
то есть множитель $2^{l-1}$ сокращается с размером подвыборки: на каждом уровне
обрабатывается вся выборка целиком. Суммируя по уровням,
\[
  T(L) = O\Bigl(d\,n\sum_{l=1}^{L}\bigl(\log n_l + c_l\bigr)\Bigr),
  \qquad
  \sum_{l=1}^{L}\log n_l = \sum_{l=1}^{L}\bigl(\log n - (l-1)\bigr)
  = L\log n - \frac{L(L-1)}{2}.
\]
Для полного дерева ($L = \log_2 n + 1$) первая сумма равна
$\tfrac12\log_2 n(\log_2 n + 1)$, а вторая приближённо
$c\log_2(n/c) + 2c$. Итого
\[
  T = O\Bigl(d\,n\,\bigl(\log^2 n + c\log n\bigr)\Bigr).
\]

\emph{(ж) Числовая оценка.} При $n = 2^{10} = 1024$, $d = 6$, $c = 16$ имеем
$\log_2 n = 10$ и $L = 11$ уровней. Считаем обе суммы точно:
\[
  \sum_{l=1}^{11}\log n_l = 10 + 9 + \dots + 1 + 0 = 55,
\]
\[
  \sum_{l=1}^{11} c_l = \underbrace{16 + 16 + 16 + 16 + 16 + 16 + 16}_{l = 1..7}
  + 8 + 4 + 2 + 1 = 112 + 15 = 127 .
\]
Число операций
\[
  T \approx d\,n\,(55 + 127) = 6 \cdot 1024 \cdot 182 \approx 1.12\cdot10^{6},
\]
и время обучения
\[
  1.12\cdot10^{6}\cdot 10^{-10}\ \text{с} \approx 1.1\cdot10^{-4}\ \text{с}.
\]

Прогнозирование устроено гораздо дешевле. Объект спускается от корня к листу,
на каждом уровне выполняется одно сравнение, глубина равна $L - 1 = 10$:
\[
  10 \cdot 10^{-10}\ \text{с} = 10^{-9}\ \text{с на объект},
\]
для всей выборки из 1024 объектов около $10^{-6}$ с. Обучение дороже
прогнозирования примерно в сто раз, и эта асимметрия -- общее свойство деревьев
решений. $\blacksquare$

\zad{6}
\dano Разбиение выборки в $i$-м узле:
\[
  \abs{S_i} = 100,\ N_i^{(1)} = 90,\ N_i^{(2)} = 10;
  \qquad
  \abs{S_{i1}} = 70,\ N_{i1}^{(1)} = 70,\ N_{i1}^{(2)} = 0;
\]
\[
  \abs{S_{i2}} = 30,\ N_{i2}^{(1)} = 20,\ N_{i2}^{(2)} = 10 .
\]

\treb Показать, что если нет разбиений, дающих большинство элементов класса
$C_2$ хотя бы в одном потомке, то неоднородность ошибки классификации равна
$0.1$ для всех возможных разбиений $S_i$, и потому не позволяет выбрать
наилучшее. Выяснить, позволяют ли это энтропия и коэффициент Джини.

\resh

\emph{Ошибка классификации.} Мера $\mathcal{I}_3 = 1 - \max_k \PP_k$ -- это доля
элементов, которые будут отнесены неверно, если весь узел пометить его
преобладающим классом. Если класс $C_1$ преобладает в обоих потомках, то
неверно классифицированными в каждом потомке оказываются ровно все элементы
$C_2$, попавшие туда. Значит абсолютное число ошибок равно общему числу элементов
$C_2$ в узле, то есть 10, независимо от разбиения. Отсюда средневзвешенная
неоднородность
\[
  \sum_j \frac{N_{ij}}{N_i}\,\mathcal{I}_3(S_{ij})
  = \frac{\text{число ошибок}}{N_i} = \frac{10}{100} = 0.1
\]
для любого такого разбиения. Неоднородность самого узла тоже
$\mathcal{I}_3(S_i) = 1 - 0.9 = 0.1$, поэтому
\[
  \Delta\mathcal{I}_3 = 0.1 - 0.1 = 0
\]
всегда. Все разбиения выглядят одинаково хорошими, выбрать среди них по этому
критерию нельзя. Проверим на данных задачи:
$\mathcal{I}_3(S_{i1}) = 1 - 70/70 = 0$,
$\mathcal{I}_3(S_{i2}) = 1 - 20/30 = 1/3$, и
$0.7\cdot 0 + 0.3\cdot\frac13 = 0.1$.

\emph{Энтропия.} Считаем в битах.
\[
  H(S_i) = -0.9\log 0.9 - 0.1\log 0.1 = 0.1368 + 0.3322 = 0.4690 .
\]
Потомки: $S_{i1}$ однороден, $H = 0$. Для $S_{i2}$ доли $2/3$ и $1/3$:
\[
  H(S_{i2}) = -\tfrac23\log\tfrac23 - \tfrac13\log\tfrac13
            = 0.3900 + 0.5283 = 0.9183 .
\]
Средневзвешенная: $0.7\cdot 0 + 0.3\cdot 0.9183 = 0.2755$, откуда
\[
  \Delta H = 0.4690 - 0.2755 = 0.1935\ \text{бита} > 0 .
\]

\emph{Джини.}
\[
  \mathcal{I}_2(S_i) = 1 - (0.9^2 + 0.1^2) = 0.18,
  \qquad
  \mathcal{I}_2(S_{i1}) = 0,
  \qquad
  \mathcal{I}_2(S_{i2}) = 1 - \Bigl(\tfrac49 + \tfrac19\Bigr) = \tfrac49 = 0.4444 .
\]
Средневзвешенная: $0.7\cdot 0 + 0.3\cdot 0.4444 = 0.1333$, откуда
\[
  \Delta\mathcal{I}_2 = 0.18 - 0.1333 = 0.0467 > 0 .
\]

Обе меры дают строго положительный прирост и, в отличие от ошибки
классификации, различают разбиения. Причина в том, что $\mathcal{I}_3$ кусочно
линейна по $\PP_k$ и зависит только от преобладающего класса, тогда как
энтропия и Джини строго вогнуты и реагируют на любое изменение долей, в том
числе на появление полностью чистого потомка $S_{i1}$. $\blacksquare$

\zad{7}
\dano Двоичное дерево решений, мера неоднородности -- энтропия.

\treb (а) Показать, что уменьшение энтропии, которое даёт один да/нет
классификатор, не больше одного бита; обобщить на произвольное число потомков.
(б) Объяснить, почему возможны ситуации, когда неоднородность в узле меньше, чем
в его потомке, и привести пример.

\resh

\emph{(а)} По \eqref{eq:mi} уменьшение энтропии есть взаимная информация
$\Delta H = I(Y; Z)$, где $Z$ -- ответ классификатора. Взаимная информация
симметрична и не превосходит энтропии каждого из аргументов:
\[
  I(Y;Z) = H(Z) - H(Z \mid Y) \le H(Z),
\]
поскольку условная энтропия неотрицательна. Для да/нет классификатора $Z$
принимает два значения, а энтропия величины с двумя исходами не больше
$\log_2 2 = 1$ бита. Значит
\[
  \Delta H \le 1\ \text{бит}. \qquad\blacksquare
\]
Равенство достигается, когда ответы делят выборку пополам и при этом полностью
определяют класс.

Для классификатора с $B$ ответами тем же рассуждением $H(Z) \le \log_2 B$, то
есть
\[
  \Delta H \le \log_2 B\ \text{бит}.
\]
Отсюда, кстати, видна причина склонности прироста информации к многозначным
атрибутам: чем больше $B$, тем больший прирост в принципе возможен, поэтому на
практике используют нормированный прирост.

\emph{(б)} Уменьшается не неоднородность каждого потомка, а её
\emph{средневзвешенное} значение, см. \eqref{eq:gain}. Отдельный потомок вполне
может оказаться неоднороднее родителя, если это компенсируется вторым, более
чистым и более крупным потомком.

Пример. Узел содержит 100 элементов: 90 класса $C_1$ и 10 класса $C_2$,
\[
  H(S) = -0.9\log 0.9 - 0.1\log 0.1 = 0.4690\ \text{бита}.
\]
Разбиение: $S_1$ -- 80 элементов, все из $C_1$; $S_2$ -- 20 элементов, 10 из
$C_1$ и 10 из $C_2$. Тогда
\[
  H(S_1) = 0, \qquad H(S_2) = 1\ \text{бит},
\]
и потомок $S_2$ неоднороднее родителя: $1 > 0.4690$. Тем не менее
\[
  \Delta H = 0.4690 - \Bigl(\frac{80}{100}\cdot 0 + \frac{20}{100}\cdot 1\Bigr)
           = 0.4690 - 0.2 = 0.2690 > 0,
\]
то есть разбиение полезно. Содержательно: вопрос выделил крупную чистую группу,
а всю неопределённость собрал в меньшую по размеру, где её и будут разбирать
дальнейшие вопросы. $\blacksquare$

\zad{8}
\dano Метод суррогатных (вспомогательных) расщеплений для случая, когда у
объекта не задано значение атрибута.

\treb Показать, что метод использует предположение, что признак, значение
которого не задано, наиболее информативен.

\resh Разберём, что именно делает метод, и какой выбор при этом не делается.

Основное разбиение узла выбирается как максимизирующее \eqref{eq:gain} по всем
атрибутам:
\[
  s = \arg\max_{i}\ \Delta\mathcal{I}(x_i).
\]
То есть уже самим построением атрибут $x_s$ объявлен самым информативным в этом
узле среди всех атрибутов.

Пусть теперь у объекта значение $x_s$ не задано. Метод суррогатных расщеплений
ищет атрибут $x_r$, разбиение по которому максимально \emph{совпадает} с
разбиением по $x_s$, и применяет его. Ключевое здесь -- критерий выбора $x_r$:
максимизируется не прирост информации $\Delta\mathcal{I}(x_r)$, а согласие с
основным разбиением
\[
  r = \arg\max_{i \ne s}\ \PP\bigl(\text{разбиение по } x_i
  \text{ совпадает с разбиением по } x_s\bigr).
\]
Иными словами, метод не переспрашивает, какой вопрос лучше задать по имеющимся
данным, а пытается \emph{угадать ответ на недоступный вопрос}. Разбиение по
$x_s$ выступает эталоном, к которому подгоняется суррогат.

Такая постановка осмысленна ровно при одном предположении: ответ на вопрос про
$x_s$ ценнее любого другого доступного ответа, поэтому приближение к нему лучше,
чем точный ответ на другой вопрос. Если бы это было не так, разумнее было бы
среди атрибутов с заданными значениями выбрать тот, что даёт наибольший
$\Delta\mathcal{I}$, и разбить по нему -- получилось бы разбиение, оптимальное
по наблюдаемым данным, но не имитирующее $x_s$.

Значит метод суррогатных расщеплений опирается на предположение о наибольшей
информативности отсутствующего признака. Заметим, что предположение это не
безобидно: атрибут выбран лучшим по обучающей выборке, где его значения были
заданы, и переносится на объекты, у которых он не задан, хотя сама пропущенность
значения может быть связана с классом. $\blacksquare$

\zad{9}
\dano Задача классификации с двумя классами и двоичными атрибутами. Обозначим
атрибуты слева направо $a$, $b$, $c$, $d$. Выборка:
\[
\begin{array}{cc}
 w_1 & w_2 \\
 0110 & 1011 \\
 1010 & 0000 \\
 0011 & 0100 \\
 1111 & 1110
\end{array}
\]

\treb (а) Построить дерево решений, используя энтропию в качестве функции
неоднородности. (б) Выразить каждый класс простейшим логическим выражением, то
есть выражением с наименьшим количеством операций OR и AND.

\resh

\emph{(а) Построение дерева.} Выборка сбалансирована: по 4 элемента в классе,
$H = 1$ бит. Считаем прирост информации для каждого атрибута в корне.

Атрибут $a$: при $a = 0$ это $\{0110,\,0011,\,0000,\,0100\}$ -- два элемента
$w_1$ и два $w_2$; при $a = 1$ то же самое. Обе половины имеют $H = 1$, значит
$\Delta H = 0$. Атрибут $b$ даёт ту же картину, $\Delta H = 0$.

Атрибут $c$: при $c = 0$ остаются $\{0000,\,0100\}$ -- оба из $w_2$, $H = 0$;
при $c = 1$ остаются шесть элементов, четыре из $w_1$ и два из $w_2$,
\[
  H = -\tfrac46\log\tfrac46 - \tfrac26\log\tfrac26 = 0.9183 .
\]
Средневзвешенная $\tfrac28\cdot 0 + \tfrac68\cdot 0.9183 = 0.6887$, откуда
$\Delta H = 1 - 0.6887 = 0.3113$ бита.

Атрибут $d$: при $d = 0$ пять элементов (два $w_1$, три $w_2$), $H = 0.9710$;
при $d = 1$ три элемента (два $w_1$, один $w_2$), $H = 0.9183$.
Средневзвешенная $0.9513$, откуда $\Delta H = 0.0487$ бита.

Максимум даёт $c$, он и становится корнем.

Ветвь $c = 0$: остались $\{0000,\,0100\}$, оба из $w_2$. Лист с меткой $w_2$.

Ветвь $c = 1$: остались $\{0110,\,1010,\,0011,\,1111\}$ из $w_1$ и
$\{1011,\,1110\}$ из $w_2$, $H = 0.9183$. Считаем прирост по оставшимся
атрибутам. Атрибут $a$: при $a = 0$ это $\{0110,\,0011\}$, оба $w_1$, $H = 0$;
при $a = 1$ это $\{1010,\,1111,\,1011,\,1110\}$, два на два, $H = 1$.
Средневзвешенная $\tfrac26\cdot 0 + \tfrac46\cdot 1 = 0.6667$, прирост
$0.9183 - 0.6667 = 0.2516$ бита. Атрибуты $b$ и $d$ делят эти шесть элементов на
части с составом $2:1$ каждая, обе с $H = 0.9183$, поэтому прирост нулевой.
Выбираем $a$.

Ветвь $c = 1$, $a = 0$: остались $\{0110,\,0011\}$, оба из $w_1$. Лист $w_1$.

Ветвь $c = 1$, $a = 1$: остались $\{1010,\,1111\}$ из $w_1$ и
$\{1011,\,1110\}$ из $w_2$. Здесь ни $b$, ни $d$ по отдельности не дают прироста:
любое из них делит четвёрку на две пары состава $1:1$. Это не тупик, а
характерная ситуация: класс определяется \emph{сочетанием} $b$ и $d$. Выпишем:
\[
\begin{array}{cccl}
  b & d & \text{вектор} & \text{класс}\\
  0 & 0 & 1010 & w_1\\
  0 & 1 & 1011 & w_2\\
  1 & 0 & 1110 & w_2\\
  1 & 1 & 1111 & w_1
\end{array}
\]
то есть $w_1$ ровно тогда, когда $b = d$ -- это задача XOR, которую один
покоординатный вопрос решить не может, а два подряд решают. Берём $b$, затем в
каждой ветви $d$.

Итоговое дерево:
\begin{center}
\begin{tabular}{ll}
\toprule
Путь & Класс \\
\midrule
$c = 0$ & $w_2$ \\
$c = 1$, $a = 0$ & $w_1$ \\
$c = 1$, $a = 1$, $b = 0$, $d = 0$ & $w_1$ \\
$c = 1$, $a = 1$, $b = 0$, $d = 1$ & $w_2$ \\
$c = 1$, $a = 1$, $b = 1$, $d = 0$ & $w_2$ \\
$c = 1$, $a = 1$, $b = 1$, $d = 1$ & $w_1$ \\
\bottomrule
\end{tabular}
\end{center}

\emph{(б) Логические выражения.} Прочитаем дерево. Элемент относится к $w_1$,
если $c = 1$ и при этом либо $a = 0$, либо $b$ совпадает с $d$:
\[
  w_1 = c \wedge \bigl(\neg a \vee (b \leftrightarrow d)\bigr).
\]
Раскроем эквивалентность через AND и OR:
\[
  \boxed{\;w_1 = c \wedge \Bigl(\neg a \;\vee\; (b \wedge d) \;\vee\; (\neg b \wedge \neg d)\Bigr)\;}
\]
Здесь три операции AND и две OR, всего пять. Для второго класса берём отрицание
и раскрываем по правилу де Моргана:
\[
  \boxed{\;w_2 = \neg c \;\vee\; \Bigl(a \wedge \bigl((b \wedge \neg d) \vee (\neg b \wedge d)\bigr)\Bigr)\;}
\]
тоже пять операций.

Проверка по всем восьми элементам:
\begin{center}
\begin{tabular}{lccccc}
\toprule
вектор & $a\,b\,c\,d$ & $c$ & $\neg a \vee (b \leftrightarrow d)$ & $w_1$ & класс \\
\midrule
$0110$ & $0\,1\,1\,0$ & 1 & $\neg a = 1$ & да & $w_1$ \\
$1010$ & $1\,0\,1\,0$ & 1 & $b = d = 0$ & да & $w_1$ \\
$0011$ & $0\,0\,1\,1$ & 1 & $\neg a = 1$ & да & $w_1$ \\
$1111$ & $1\,1\,1\,1$ & 1 & $b = d = 1$ & да & $w_1$ \\
$1011$ & $1\,0\,1\,1$ & 1 & $a = 1$, $b \ne d$ & нет & $w_2$ \\
$0000$ & $0\,0\,0\,0$ & 0 & --- & нет & $w_2$ \\
$0100$ & $0\,1\,0\,0$ & 0 & --- & нет & $w_2$ \\
$1110$ & $1\,1\,1\,0$ & 1 & $a = 1$, $b \ne d$ & нет & $w_2$ \\
\bottomrule
\end{tabular}
\end{center}
Все восемь элементов классифицированы верно. $\blacksquare$

Заметим, что на этой же выборке работает и более короткая по виду формула
$w_1 = c \wedge (a \oplus b \oplus d)$, но операция XOR в постановке не
разрешена, а её раскрытие через AND и OR требует одиннадцати операций, то есть
больше приведённых пяти.

\end{document}
